\subsection{Exploratory Data Analysis}
\label{subsec:data_eda}

\begin{sloppypar}

Before any preprocessing decision is made, each of the three raw datasets is
examined independently in \texttt{notebooks/\allowbreak 01\_eda.\allowbreak ipynb} to characterize its
class balance, missingness, numeric-feature relationships, and the strength of
association between individual features and the target. The purpose of this
stage is not to fit any model, but to surface the data-quality issues and
predictive signals that later motivate specific preprocessing choices (see
Section~\ref{subsec:data_preprocessing}) and the class-imbalance policy adopted
in the modeling stage (Section~\ref{subsec:data_description} and finding F1 of
the project's scope review).

\paragraph{Class imbalance.} All three datasets exhibit some degree of class
imbalance, but to very different extents. Dataset 1 is the most imbalanced,
with a raw positive rate of 5.64\%, giving balanced class weights of
approximately $\{0: 0.53, 1: 8.80\}$ -- a positive example is weighted roughly
16.6$\times$ more than a negative one under scikit-learn's balanced weighting
scheme \cite{pedregosa2011scikit}. Dataset 3 is imbalanced in the same
direction but less severely, with a raw positive rate of 9.42\% and balanced
weights of approximately $\{0: 0.55, 1: 5.31\}$ (about 9.6$\times$). Dataset 2
is the outlier among the three: its positive rate of 55.34\% is close to
balanced, giving weights near 1.0 in both directions ($\{0: 1.12, 1: 0.90\}$)
and requiring very little correction. This spread of imbalance regimes across
the three datasets is precisely what motivates RQ1/RQ2's design: a class-
imbalance policy that only worked well on a near-balanced dataset like Dataset
2 would not generalize to the 17:1 imbalance of Dataset 1.

\paragraph{Missingness.} Dataset 1 is the only one of the three with
meaningful missing data: 38 of its 40 columns contain at least one missing
value, led by vaccination/screening-history fields that respondents often
skip (Table~\ref{tab:d1_missing}). Datasets 2 and 3 report zero missing values
on the surface, but Dataset 2's completeness is misleading -- its
\texttt{RestingBP} and \texttt{Cholesterol} columns use an undocumented
\texttt{0} placeholder for "not measured" (1 and 172 rows respectively) that
is not distinguishable from a genuine reading without domain knowledge. This
is why the EDA stage treats "zero missing values" and "no data-quality
issues" as two separate questions rather than assuming one implies the other.

\paragraph{Numeric feature correlations.} Correlation matrices computed over
each dataset's numeric columns show no evidence of problematic
multicollinearity in any of the three datasets, though each has its own
notable pairwise relationships. In Dataset 1, \texttt{WeightInKilograms} and
\texttt{BMI} are strongly correlated ($r = 0.86$), which is expected since BMI
is mathematically derived from weight and height; \texttt{HeightInMeters} and
\texttt{WeightInKilograms} are moderately correlated ($r = 0.47$), while
\texttt{HeightInMeters} and \texttt{BMI} are nearly uncorrelated
($r = -0.03$), confirming that BMI successfully normalizes out the effect of
height. \texttt{PhysicalHealthDays} and \texttt{MentalHealthDays} are
moderately correlated ($r = 0.32$). In Dataset 2, no pair of numeric features
is strongly correlated; the strongest relationships are \texttt{Age}
$\leftrightarrow$ \texttt{MaxHR} ($r = -0.38$) and \texttt{Age}
$\leftrightarrow$ \texttt{Oldpeak} ($r = 0.26$), with \texttt{Cholesterol}
weakly negatively correlated with \texttt{FastingBS} ($r = -0.26$), plausibly
diluted by the placeholder-zero \texttt{Cholesterol} values noted above. In
Dataset 3, \texttt{PhysHlth} and \texttt{GenHlth} are the most correlated pair
($r = 0.52$), \texttt{MentHlth} and \texttt{PhysHlth} are moderately
correlated ($r = 0.35$), and \texttt{Education} and \texttt{Income} show the
strongest socioeconomic correlation overall ($r = 0.45$). In all three cases,
these pairwise relationships are informative but well below the
$0.9$ correlation threshold later used for feature reduction
(Section~\ref{subsec:data_preprocessing}), so they surface as candidates for
the modeling stage rather than as columns to drop outright at the EDA stage.

\paragraph{Outlier detection (Dataset 2 only).} Dataset 2 is the only dataset
whose numeric columns are genuine continuous clinical measurements rather than
mostly binary or discrete survey answers, which is why IQR (Tukey-fence)
outlier detection is examined here and later applied only to Dataset 2 during
preprocessing (finding F5). Computed on the full raw dataset for exploratory
purposes, \texttt{Age} has no values outside its fence; \texttt{RestingBP}
(bounds 90--170) flags 27 potential outliers; \texttt{Cholesterol} (bounds
105.6--376.6, computed after treating its zero-placeholders as missing) flags
23; \texttt{MaxHR} (bounds 66--210) flags 2; and \texttt{Oldpeak} (bounds
$-2.25$--3.75) flags 16, consistent with its long right tail. These
full-dataset bounds are for exploratory characterization only; the bounds
actually used to clip Dataset 2's train/test splits in
Section~\ref{subsec:data_preprocessing} are refit on the training split alone
to avoid leaking test-set extremes into the fences.

\paragraph{Target rate by category.} Categorical features associated with
self-reported general health status are, across every dataset that has one,
by far the strongest single predictors observed during EDA. In Dataset 1, the
heart-attack rate rises monotonically with self-reported \texttt{GeneralHealth}
from 1.53\% ("Excellent") to 2.84\% ("Very good"), 5.78\% ("Good"), 11.99\%
("Fair"), and 21.08\% ("Poor") -- a roughly 14-fold spread across five
categories, in contrast to BMI, whose distribution is nearly identical between
positive and negative cases (median around 27--28 in both groups) and is
therefore a comparatively weak discriminator on its own. In Dataset 2, both
categorical predictors examined are strong: patients with asymptomatic chest
pain (\texttt{ASY}) have a heart-disease rate of 79.03\%, far higher than
\texttt{ATA} (13.87\%), \texttt{NAP} (35.47\%), or \texttt{TA} (43.48\%) --
counterintuitively, the absence of typical chest-pain symptoms is the
strongest single warning sign in this dataset -- and a \texttt{Flat} (82.83\%)
or \texttt{Down} (77.78\%) ST-segment slope is associated with a much higher
rate than an \texttt{Up}-sloping segment (19.75\%). In Dataset 3, the same
\texttt{GenHlth} pattern recurs almost identically to Dataset 1 despite the two
datasets being independent BRFSS survey years: the rate rises from 2.24\%
("Excellent") to 4.63\%, 10.46\%, 21.31\%, and 34.00\% ("Poor"); \texttt{HighBP}
is also strongly associated, with 16.47\% of respondents with high blood
pressure reporting heart disease/attack versus only 4.12\% of those without.

Taken together, these EDA findings motivate three concrete preprocessing
decisions detailed in Section~\ref{subsec:data_preprocessing}: a shared
class-imbalance policy applied at model-fit time rather than by resampling
(F1, since the degree of imbalance varies too much across datasets for a
single resampling ratio to be appropriate); IQR-based outlier clipping scoped
to Dataset 2 only (F5, since Datasets 1 and 3 are mostly binary/discrete
survey data where a genuinely high value is itself a disease signal rather
than noise); and dataset-specific data-quality handling -- zero-as-missing
recoding for Dataset 2 and duplicate-row removal for Dataset 3 -- rather than
a single one-size-fits-all cleaning rule.
\end{sloppypar}
